\section{Frontend: Compilación, Optimización Final y Distribución}
\label{sec:s5_frontend}

La compilación final optimiza el \textit{bundle} y lo integra en el artefacto monolítico. La
Figura~\ref{fig:s5_build} ilustra el \textit{pipeline} de empaquetado FE+BE.\par

\begin{figure}[!ht]
    \centering
    \begin{tikzpicture}[
        node distance=0.55cm,
        s/.style={draw=colorPrimario, fill=headerTabla, thick, rounded corners=2pt,
            text width=2.3cm, align=center, minimum height=1.0cm, font=\sffamily\scriptsize},
        rel/.style={-{Latex[length=2mm]}, colorPrimario, font=\sffamily\tiny}
    ]
        \node[s] (a) {1. tsc + vite\\build};
        \node[s, right=of a] (b) {2. dist/\\(minificado)};
        \node[s, right=of b] (c) {3. copia a\\static/};
        \node[s, right=of c] (d) {4. JAR\\monolítico};
        \draw[rel] (a)--(b); \draw[rel] (b)--(c); \draw[rel] (c)--(d);
    \end{tikzpicture}
    \caption{Pipeline de compilación y empaquetado del frontend en el JAR.}
    \label{fig:s5_build}
\end{figure}

\subsection{Compilación Optimizada y Code Splitting}
\label{ssec:s5_fe_build}
El proceso de compilación (\comando{tsc \&\& vite build}) genera \comando{dist/} con
minificación de assets, remoción de dependencias muertas y división del código
(\textit{Code Splitting}); el \textit{chunk} \comando{vendor} separa \comando{react},
\comando{react-dom} y \comando{react-router-dom} para mejorar el cacheo. La adaptabilidad final
incluye i18n (es/en/pt) y modo claro/oscuro.

\begin{itemize}[noitemsep]
    \item \textbf{Minificación:} reducción de JS/CSS y \textit{tree-shaking}.
    \item \textbf{Code splitting:} \textit{chunk} \comando{vendor} para librerías estables.
    \item \textbf{Adaptabilidad:} cambio de idioma y de tema sin recarga.
\end{itemize}

\subsection{Distribución de Artefactos Estáticos}
\label{ssec:s5_fe_cdn}
Los artefactos estáticos se sirven desde el JAR de Spring Boot (despliegue monolítico) y pueden
distribuirse a través de una red de distribución de contenido (CDN) para garantizar mínima
latencia. El \comando{base} de Vite es \comando{/}.

\begin{NotaTecnica}
El despliegue monolítico copia \comando{dist/} a \comando{ethoshubB/src/main/resources/static/};
el \comando{SpaFallbackFilter} sirve \comando{index.html} para las rutas de la SPA, de modo que
la navegación del cliente no produce errores 404 en recargas profundas.
\end{NotaTecnica}

\subsection{Optimización del Bundle e Internacionalización Final}
\label{ssec:s5_fe_bundle}
La compilación de producción aplica las optimizaciones de la tabla~\ref{tab:s5_bundle} y completa
la adaptabilidad (i18n y tema).

\begin{table}[!ht]
    \centering
    \renewcommand{\arraystretch}{1.3}
    \fontsize{9}{10.8}\selectfont\sffamily
    \begin{tabularx}{\textwidth}{|l|X|}
        \hline
        \rowcolor{colorPrimario}
        \textcolor{white}{\textbf{Optimización}} & \textcolor{white}{\textbf{Efecto}} \\ \hline
        Minificación & Reducción de JS/CSS y \textit{tree-shaking}. \\ \hline
        \rowcolor{grisTecnico}
        \textit{Vendor chunk} & Separa \comando{react}, \comando{react-dom}, \comando{react-router-dom}. \\ \hline
        \textit{Code splitting} & Carga diferida por ruta. \\ \hline
        \rowcolor{grisTecnico}
        i18n (es/en/pt) & Cambio de idioma sin recarga; \textit{fallback} a \comando{es}. \\ \hline
        Tema claro/oscuro & Variables de tema persistentes. \\ \hline
    \end{tabularx}
    \caption{Optimizaciones del bundle y adaptabilidad final (Sprint 5).}
    \label{tab:s5_bundle}
\end{table}

El \textit{bundle} resultante se copia a \comando{static/} del backend para el despliegue
monolítico; el cacheo del \textit{vendor chunk} reduce el tiempo de carga en visitas recurrentes,
y la distribución vía CDN minimiza la latencia para usuarios geográficamente dispersos.

\clearpage
